\psection{protocol}{0.9}{I1--I4}{core}{Protocol}
% Section V. Sources checked on 16 Sep: VF config/deployments.yaml,
% config/foundry.yaml, config/experiments/E1_panel.yaml, E2_seeds.yaml,
% E6_scorecard.yaml; VF docs/KNOWN_LIMITS.md sections 2, 3 and 11; VF
% analysis/measurability.py and causes.py; RAIP VERA_questionnaire_avant_apres.md.
% The APSEC reading study, its counts and its ethics sentence are not reused.

% --- P1: Arm L (design established; versions and dates pending E0/E1) -------
\textbf{LLM arm: language models on Microsoft Foundry.}
The arm evaluates eight deployments in three tiers: five frontier models from Anthropic, OpenAI and Mistral, an open-weight 70B model, and two small models (8B and 3B).
All but the OpenAI model, served by Azure OpenAI, use serverless endpoints; the 70B model may move to managed compute \tbd{E1}{foundry}{27 Sep}{serving mode the 70B deployment actually used (serverless or managed compute)}.
Each deployment is scored on the nine output-dependent requirements (\req{01}, \req{02}, \req{06}--\req{12}); the data requirements \req{03}--\req{05} are scored once, on a synthetic banking corpus.
Runs use temperature 0, at most 256 output tokens, 150 items per benchmark, seed 42 and 500 bootstrap resamples~\cite{efron1993}, and four deployments receive two further seeds.
A fixed judge, Llama~3.3 70B served on Azure, scores refusals and attack success.
The client substitution and runtime patches this required are reported with the portability results\armnotwithdrawn{ (\Cref{sec:portability})}.
Model versions are \tbd{E0}{foundry}{18 Sep}{model version of each deployment, from the E0 probe with write enabled} and run dates \tbd{E1}{foundry}{27 Sep}{run dates of the eight panel runs}.

% --- P2: Arm M (pending N1; depends on the 30 Sep decision) -----------------
\textbf{Tabular arm: credit-scoring classifiers.}
\armswitch{%
\parplan[40]{The tabular arm trains reference classifiers on a synthetic credit corpus with a planted disparity and proxy, on UCI Adult and on South German Credit, and gives VERA only their predictions.}{N1 (G1, G9). A predictions file with true label, predicted label, score and protected attribute keeps the engine independent of the model; only counterfactual checks would need the model artefact (D3). Cite becker1996adult, hofmann1994statlog with gromping2019south (corrected South German coding), ding2021retiring. Protected attribute: D6 (22 Sep). Reference models and split PENDING. Scope: fairness, calibration, data adequacy, privacy (D2). The synthetic manifest gives ground truth for the planted effects. Mirror Arm L: model versions, seed, data sizes. Runs due 4 Oct.}%
}{%
\parplan[25]{The tabular arm trains one reference classifier on the synthetic credit corpus, with a planted disparity, and gives VERA only its predictions.}{State that the public datasets are specified but not run. N1 reduced scope: one dataset, fairness criteria and calibration only; list what ran and mark the rest specified, not run. Protected attribute D6. Point estimates unless G4 lands. Runs due 4 Oct.}%
}{%
The tabular arm was designed (\Cref{sec:vera}) but not run, so it contributes no measured result.%
}

% --- P3: criterion sensitivity (procedure established; C and thresholds D5) -
\textbf{Criterion sensitivity.}
Each admissible fairness criterion $c \in \mathcal{C}$ is a separate benchmark under \req{11} with a score $m_c \in [0,1]$, higher meaning fairer, so adopting one criterion is a one-hot weighting and mixing criteria is any other weighting.
A verdict is criterion-robust if all criteria in $\mathcal{C}$ and all their mixtures agree on it; \Cref{sec:fairness} gives the test, including per-criterion thresholds $\tau_c$.
\legalrisk{} fixes $\mathcal{C}$ and the threshold rule \decision{D5}{legal}{30 Sep}{Admissible criteria; common band scale or per-criterion thresholds; default criterion for credit decisions; a rationale that can be quoted.}.
The test covers the LLM arm, whose language models classify credit records rendered as text and feed the same metrics module \decision{D7}{core}{22 Sep}{LLM fairness instrument: record classification by the language models (N2), or the existing R10 and R11 probes with caveats.}\armswitch{, and every model--dataset pair of the tabular arm \tbd{N3}{core}{9 Oct}{criterion-sensitivity analysis on both arms, with the per-criterion threshold variant; row-level bootstrap intervals (G4) or point estimates declared as a threat}}{, and point estimates for the synthetic dataset of the tabular arm \tbd{N3}{core}{9 Oct}{criterion-sensitivity analysis on the LLM arm and on tabular point estimates}}{ \tbd{N3}{core}{9 Oct}{criterion-sensitivity analysis on the LLM record-classification runs (N2) only}}.

% --- P4: measurability (procedure established; E0/N4 values pending) --------
\textbf{Measurability.}
Each deployment--requirement cell is native (the catalogued harness ran without a fallback flag), fallback (a degraded path, with its logged reason) or not measurable; native thus means unflagged, not valid.
Each non-native cell gets a cause class from its logged reason: serving (the endpoint lacks a needed feature, such as log-probabilities), set-up (a tool is missing or unusable in the campaign set-up), by design (no native harness, or not applicable) or unknown.
Unrecognised reasons never count as serving, and only serving-class cells support serving-mode claims.
A probe records log-probability and seed support per deployment \tbd{E0}{foundry}{18 Sep}{log-probability and seed support per deployment}.
With \legalrisk{}, we classify each of the twelve measurable requirements as transferable to tabular classifiers or not \tbd{N4}{core,legal}{30 Sep}{family-transfer classification of the twelve measurable requirements, one reason each (D2)}; \armswitch{Tabular-arm runs then fill the tabular column of the matrix}{Tabular-arm runs fill the tabular column of the matrix where they exist, and other cells are marked not run}{without Tabular-arm runs, this analytical classification is the tabular column of the matrix}.

% --- P5: decision study (design established from the questionnaire) --------
\textbf{Decision study.}
In a within-subject design~\cite{wohlin2012experimentation}, participants choose one of three anonymised models for a customer-facing retail-banking assistant, in two phases.
Phase one shows what a reviewer usually has (size, licence, a capability score, cost, latency, vendor claim); participants choose, justify, rate confidence from 1 to 5 and name missing information (Q1--Q4).
Phase two, presented as a second source of information without naming the tool, adds requirement scores in three colour bands and marks verdicts that another defensible weighting would change \decision{D11}{core,foundry}{18 Sep}{Which scorecard was sent on 14 Sep (the older Mistral trio or the E6 Foundry trio); path and timestamp of the freeze record.}.
They choose again, explain any change, rate confidence and tick the criteria that weighed (Q5--Q8), and cannot return to phase one.
Closing items cover declared review effort, evidence sources, usefulness and role (Q9--Q16).
\armnotwithdrawn{\legalrisk{} participants also answer a tabular credit-scoring variant \decision{D14}{legal,core}{25 Sep}{Build or drop the decision instrument for the tabular arm.}.}

The measures are M1, the switch rate, with a Wilson interval; M2, agreement with the choice fixed before collection in a dated position note \decision{D12}{core}{18 Sep}{Does a position note dated before diffusion exist? Otherwise reframe or drop M2.}, by McNemar's test~\cite{mcnemar1947}; M3, decision time, which is not analysed because the form timestamps whole phases, not the choice; M4, confidence, by the Wilcoxon signed-rank test~\cite{wilcoxon1945} with the rank-biserial correlation~\cite{kerby2014}; M5, the count of distinct criteria cited in Q2 and Q6, coded twice \decision{D35}{core,dal}{5 Oct}{Name the two M5 coders (both must read French, or declare the translation of Q2 and Q6 before coding) and fix the M5 rule for respondents who kept their choice.} with Krippendorff's~$\alpha$ for agreement~\cite{hayes2007answering} and compared by the same test; M6, declared review effort; and M7, perceived usefulness.
Tests run on the pooled sample, and Holm's procedure~\cite{holm1979} adjusts M2, M4 and M5 as one family; population contrasts are descriptive and exploratory.
The study targets \aicc{} and \legalrisk{} \decision{D13}{core,aicc,legal}{18 Sep}{How respondents are assigned to \aicc{} or \legalrisk{} (tracked channel or an added question).}, with at least 30 participants and a floor, fixed between 15 and 20 at collection close before the export is opened, below which we report descriptive statistics only \tbd{N6}{core}{28 Sep}{collection window, invitations sent and responses per population}.

% --- P6: case study (pending N8 and D16) ------------------------------------
\textbf{Case study.}
\parplan[10]{We follow one \aicc{} use case from submission to a documented review decision, as a single case whose embedded units of analysis are the requirements (\Cref{sec:casestudy})~\cite{runeson2009guidelines}.}{N8 and D16 (case choice and permission 25 Sep; go or no-go on a real review 2 Oct). Section VII owns the design and the data sources; keep one sentence here. Never claim that a run fed a release review before it has. Degraded outcome: a labelled retrospective of a past decision.}

% --- P7: ethics (approval clauses pending D15) ------------------------------
\textbf{Ethics.}
The form collects no names, responses are analysed in aggregate, and no run in either arm uses customer or employee data.
Collection requires approval by the bank's Responsible AI and Risk and Compliance functions, managers' agreement and informed consent \decision{D15}{legal}{25 Sep}{Compliance and legal sign-off, manager agreement, consent text, form tool and English version; must precede further collection.}.
We declare compliance with the ACM policy on human participants at submission; the study is not pre-registered.
